\documentclass[12pt]{article}

\usepackage{sbc-template}
\usepackage{graphicx,url}
\usepackage[utf8]{inputenc}
\usepackage[brazil]{babel}
\usepackage{helvet}
\usepackage{caption}
\usepackage{booktabs}
\usepackage{multirow}



\usepackage{times}
\usepackage[T1]{fontenc}
\usepackage{url}
\usepackage{booktabs}
\usepackage{multirow}
\usepackage{graphicx}
\pagenumbering{gobble}
\setlength{\parindent}{1.27cm}
\makeatletter
\renewcommand\@biblabel[1]{}
\makeatother

\DeclareCaptionFont{sbcfont}{\sffamily\bfseries\fontsize{10}{12}\selectfont}
\captionsetup{font=sbcfont, aboveskip=6pt, belowskip=6pt}

     
\sloppy

\title{Radiomics applied to skin lesion classification in HFUS images}

\author{Lucas G. Guedes\inst{1}, Isabela R. V. da Silva\inst{2}, Thatiane A. P. Alva\inst{3}, Viviane R. Botelho\inst{4}}


\address{Departamento de Ciências Exatas e Sociais Aplicadas\\
  Universidade Federal de Ciências da Saúde de Porto Alegre\\
  R. Sarmento Leite, 245 -- 90050-170 -- Porto Alegre -- RS -- Brazil
  \email{\{lucas.guedes, isabelarv, thatiane, vivianerb\}@ufcspa.edu.br}
}

\begin{document} 

\maketitle

\begin{abstract}
This study evaluates the performance of classical machine learning classifiers trained with radiomic features extracted from high-resolution ultrasound images, to classify cutaneous lesions according to their malignancy. First-order and texture features were extracted, with selection performed via ANOVA F-score and Mutual Information, and 23 classifiers were evaluated in three datasets. The best model, a Gaussian SVM, achieved an accuracy of $95.0\%$, sensitivity of $100\%$ and a performance equivalent to deep learning architectures reported in the literature. The results ndicate that radiomic features are sufficient to match or exceed the performance of dedicated CNNs, with the advantages of lower computational cost and greater interpretability. 
\end{abstract}


\section{Intruduction}

Skin cancer is the most frequently diagnosed neoplasm in the world and also most complex in humans \cite{Juan2023, Zhu2024}. When detected in early stages, the overall survival rate can reach $99\%$, but this effectiveness decrease as the disease progress \cite{Goyal2020}. To make this possible, dermatologists evaluate lesions through different analysis methods as dermatoscopy, confocal microscopy and other images techniques; however, this images are limited in terms of depth of lesion \cite{Xiao2025, Zhu2024}. High-frequency ultrasound helps overcome this limitation by providing high-resolution images from the epidermis to deep fascia, showing potential for differentiating benign and malignant skin lesions \cite{daSilva2026, Czajkowska2021-1}. However, their interpretation remains subjective and dependent on specialist experience, which can lead to diagnostic errors and delays in
reaching a conclusion \cite{Qin2025}

Given this scenario, Machine Learning (ML) and Deep Learning (DL) techniques have emerged as promising tools to assist medical image diagnosis \cite{daSilva2026}. Algorithms based on convolutional neural networks have been widely applied to various medical images modalities, demonstrating performance comparable to human specialists in classification and detection tasks \cite{LaverdeSaad2021}. Despite these promising results, deep learning models present a relevant limitation for clinical application: their predictions are often difficult to interpret, since the feature extraction process occurs implicitly within the neural network \cite{Qin2025}. Furthmore, effictive neural network training requires a hing computational and data cost \cite{LaverdeSaad2021}. 

In this context, radiomics emerge as a extraction process of a large number of features related to morphological and textural information from medical images, making quantitative information about lesions accessible that is not perceptible through direct visual assessment \cite{Lambin2017}. These extracted characteristics, which include intensity, shape, texture, and wavelet transform values, capture characteristics of the lesion's phenotype that are distinct from or complementary to information obtained by other methods \cite{Lambin2017}.

Despite growing interest in radiomics for medical imaging, few studies have applied it to ultrasound-based skin tumor classification, with most relying on DL instead \cite{Surkov2024, Bibi2026, Czajkowska2021-2}. Given this gap, the present work aims to evaluate the performance of several classical machine learning classifiers trained with
radiomic features extracted from HFUS images to assess the malignancy of
cutaneous lesions, comparing the results obtained with those reported by Silva
et al. (2026)

\section{Materials and Methods}
\subsection{Dataset and Feature Extraction}
The dataset used was produced by Alfageme et al. \cite{Alfageme2022}, comprising 240 $\times$ 240 px PNG HFUS images (BW and Doppler modes) of cutaneous lesions, labeled as benign (65.7\%) or malignant (34.3\%). Images were stratified into 70\% train / 30\% test subsets to preserve class proportions. Prior to extraction, dataset inconsistencies documented by Silva et al. (2026) were corrected to ensure direct comparability with their results.

Radiomic features were extracted using PyRadiomics, yielding 465 first-order and texture features per image, calculated on the original image and on the four sub-bands of the 2D wavelet decomposition (LL, LH, HL, HH), which capture spatially localized frequency content otherwise suppressed in the original image. Three feature datasets were created: BW, Doppler, and combined BW+Doppler.

\subsection{Experimental Design and Evaluation}
Given the high number of features, SelectKBest (scikit-learn) was used to rank feature relevance on the training set with two criteria: the ANOVA F-score (inter-/intra-class variance ratio) and Mutual Information (non-linear feature-target dependence), both computed on the training set only to avoid data leakage.

This motivated two parallel scenarios: Track A, in which standardized features were reduced via SelectKBest to the K most discriminative ones (K = 5-100) before training; and Track B, in which all standardized features were retained. The same 23 classifiers were applied to both tracks across the three datasets (BW, Doppler, BW+Doppler). Performance was estimated via stratified 5-fold cross-validation on the training set, with feature selection in Track A refit within each fold to avoid leakage between training and validation. The classifier with the highest mean cross-validation AUC per track and dataset was retrained on the full training set and evaluated once on the held-out test set, using both a 0.5 threshold and the Youden-index threshold; AUC, accuracy, sensitivity, specificity, and F1-score were computed for each.

\section{Results and Discussion}

Table \ref{tab:resultados-youden} present the best-performing per track and dataset on the held-out test set. In both tracks, the top result was obtained using the combined BW+Doppler dataset.

\begin{table}[ht]
\centering
\caption{Performance of the best classifiers per dataset in youden threshold (on the test data) - Tracks A and B}
\label{tab:resultados-youden}
\resizebox{0.8\textwidth}{!}{%
\begin{tabular}{clcccccc}
\toprule
\textbf{Track} & \textbf{Dataset} & \textbf{Model} & \textbf{Sens.} & \textbf{Spec.} & \textbf{F1} & \textbf{Acc.} & \textbf{AUC} \\
\midrule
\multirow{3}{*}{A} & BW         & Coarse Gaussian SVM    & 0.71 & 1.00 & 0.83 & 0.90 & 0.87 \\
                    & Doppler    & Subspace Discriminant  & 1.00 & 0.69 & 0.78 & 0.90 & 0.99 \\
                    & BW+Doppler & Subspace Discriminant  & 0.71 & 0.85 & 0.71 & 0.90 & 0.96 \\
\midrule
\multirow{3}{*}{B} & BW         & Medium Gaussian SVM    & 0.71 & 0.92 & 0.77 & 0.85 & 0.96 \\
                    & Doppler    & Medium Gaussian SVM    & 0.57 & 1.00 & 0.73 & 0.85 & 0.88 \\
                    & BW+Doppler & Medium Gaussian SVM    & 1.00 & 0.92 & 0.93 & 0.95 & 0.98 \\
\bottomrule
\end{tabular}}
\end{table}

The best-performing model (Medium Gaussian SVM, Track B, BW+Doppler) achieved an accuracy of 95.0\%, sensitivity of 100\%, specificity of 92.3\%, and AUC of 0.98, matching with accuracity and AUC reported by Silva et al. (2026) for their best CNN (HFUS-Doppler: 95.0\% accuracy, AUC 0.98). While the SVM model also expressively surpass the EfficientNet B4 (accuracy: $77.1\%$, sensitivity: $73.3\%$, specificity: $80.0\%$) model reported by Laverde-Saad et al. \cite{LaverdeSaad2021}. 

The classifier's performance compared to a deep neural network trained directly on image pixels suggests that the textural and morphological information captured by radiomics is sufficient to match dedicated CNNs in discriminative capacity, with the advantage of maintaining greater model interpretability and substantially lower computational cost. The best classifier trained in approximately 51ms on a conventional CPU (Intel Core i5-7267U, 8~GB RAM), without GPU usage

A limitation of this study is the small test set (n=20: 13 malignant, 7 benign), which limits the stability of point estimates; results should therefore be regarded as preliminary evidence requiring confirmation on larger datasets.

This work shows that classical machine learning classifiers trained on radiomic features extracted from HFUS images can match or exceed CNN-based approaches for benign/malignant skin lesion classification, while requiring substantially less computational infrastructure and offering greater interpretability. Future work should validate these findings on larger, multi-center datasets.

\begin{thebibliography}{99}
\setlength{\leftmargin}{0.5cm}
\setlength{\itemindent}{-0.5cm}
\setlength{\itemsep}{6pt}
\setlength{\parsep}{0pt}
\bibitem[Alfageme 2022]{Alfageme2022} Alfageme, F. (2022). Dermatologic ultrasound images for classification. \url{https://www.kaggle.com/datasets/alfageme/dermatologic-ultrasound-images}.

\bibitem[Bibi et al. 2026]{Bibi2026} Bibi, F., Abbas, Q. and Naqi, S. M. (2026) dual-stream deep learning model for early detection and classification of non-dermoscopic skin lesions via focal modulation and local feature integration. \textit{IEEE Access}, 14:41998--42014.

\bibitem[Czajkowska et al. 2021a]{Czajkowska2021-1} Czajkowska, J., Badura, P., Korzekwa, S., P{\l}atkowska-Szczerek, A. and S{\l}owi{\'n}ska, M. (2021). Deep learning-based high-frequency ultrasound skin image classification with multicriteria model evaluation. \textit{Sensors}, 21(17):5846.

\bibitem[Czajkowska et al. 2021b]{Czajkowska2021-2} Czajkowska, J., Badura, P., Korzekwa, S. and P{\l}atkowska-Szczerek, A. (2021). Deep learning approach to skin layers segmentation in inflammatory dermatoses. \textit{Ultrasonics}, 114:106412.

\bibitem[Goyal et al. 2020]{Goyal2020} Goyal, M., Knackstedt, T., Yan, S. and Hassanpour, S. (2020). Artificial intelligence-based image classification methods for diagnosis of skin cancer: challenges and opportunities. \textit{Computers in Biology and Medicine}, 127:104065.

\bibitem[Juan et al. 2023]{Juan2023} Juan, C. K., Su, Y. H., Wu, C. Y., et al. (2023). Deep convolutional neural network with fusion strategy for skin cancer recognition: model development and validation. \textit{Scientific Reports}, 13(1):17087.

\bibitem[Lambin et al. 2017]{Lambin2017} Lambin, P., Leijenaar, R. T., Deist, T. M., et al. (2017). Radiomics: the bridge between medical imaging and personalized medicine. \textit{Nature Reviews Clinical Oncology}, 14(12):749--762.

\bibitem[Laverde-Saad et al. 2021]{LaverdeSaad2021} Laverde-Saad, A., Jfri, A., Garc\'ia, R., et al. (2021). Discriminative deep learning based benignity/malignancy diagnosis of dermatologic ultrasound skin lesions with pretrained artificial intelligence architecture. \textit{Skin Research and Technology}, 28(1):35--39.

\bibitem[Qin et al. 2025]{Qin2025} Qin, Y., Zhang, Z., Qu, X., Liu, W., Yan, Y. and Huang, Y. (2025). A machine learning model based on high-frequency ultrasound for differentiating benign and malignant skin tumors. \textit{Medical Ultrasonography}, 27(3):284--293.

\bibitem[Silva et al. 2026]{daSilva2026} Silva, I. R. V. da, Jardim, A. G., Fa\'es, G. R. da S., Alva, T. A. P., Becker, C. D. L. and Botelho, V. R. (2026). Deep learning-based skin lesion classification: a CNN approach on high-frequency ultrasound imaging. \textit{Journal of Ultrasound in Medicine}, 45(5):925--936.

\bibitem[Surkov et al. 2024]{Surkov2024} Surkov, Y. I., Serebryakova, I. A., Kuzinova, Y. K., et al. (2024). Multimodal method for differentiating various clinical forms of basal cell carcinoma and benign neoplasms in vivo. \textit{Diagnostics}, 14(2):202.

\bibitem[Xiao et al. 2025]{Xiao2025} Xiao, C., Zhu, A., Xia, C., et al. (2025). Attention-guided learning with feature reconstruction for skin lesion diagnosis using clinical and ultrasound images. \textit{IEEE Transactions on Medical Imaging}, 44(1):543--555.

\bibitem[Zhu et al. 2024]{Zhu2024} Zhu, A. Q., Wang, Q., Shi, Y. L., et al. (2024). A deep learning fusion network trained with clinical and high-frequency ultrasound images in the multi-classification of skin diseases in comparison with dermatologists: a prospective and multicenter study. \textit{eClinicalMedicine}, 67:102391.
\end{thebibliography}


\end{document}
